\section*{Hoofdstuk \ref{ch:b3:sensory-systems} --- Zintuiglijke stelsels}

\begin{solution}{exo:b3:sensory-systems:1}
Aan de draad, niet aan de boodschap: elke zintuiglijke zenuw is een
\omterm{def:b3:sensory-systems:principles}{gemerkte lijn} die in haar eigen gebied van het brein eindigt, en wat er
langs de oogzenuw aankomt wordt als licht gelezen --- en daarom
veroorzaakt druk op de oogbol een lichtflits en een klap op het oor een
piep.
\end{solution}

\begin{solution}{exo:b3:sensory-systems:2}
In het donker houdt cyclisch GMP kationkanalen open en houdt een
gestage inwaartse ``donkerstroom'' de cel gedepolariseerd en glutamaat
afgevend. Licht activeert \omterm{def:b3:sensory-systems:retina}{rodopsine}, transducine en fosfodi-esterase,
dat cGMP afbreekt; de kanalen sluiten, de inwaartse stroom stopt, en de
cel hyperpolariseert en geeft minder af --- licht wordt door een afname
gemeld.
\end{solution}

\begin{solution}{exo:b3:sensory-systems:3}
$I = I_{0}10^{L/10}$: \qty{20}{dB}, \qty{1e-10}{W/m^2}; \qty{60}{dB},
\qty{1e-6}{W/m^2}; \qty{110}{dB}, \qty{0.1}{W/m^2}. Van het luidste tot
het zachtste: $10^{9}$.
\end{solution}

\begin{solution}{exo:b3:sensory-systems:4}
Zoet, umami en bitter via receptoren gekoppeld aan een G-eiwit (één
voor zoet, één voor umami, ongeveer vijfentwintig voor bitter); zout
via het natriumkanaal ENaC; zuur via een protonkanaal. Het meeste van
de ``\omterm{def:b3:sensory-systems:chemical}{smaak}'' is de reuk van vluchtige moleculen die vanuit de mond de
neus bereiken; met een verstopte neus blijven alleen de vijf smaken
over en lijkt voedsel flauw.
\end{solution}

\begin{solution}{exo:b3:sensory-systems:5}
Fechner: $S = (1/0.02)\ln(10\,000/100) = 50\ln 100 \approx 230$
stappen. Tellen: $1.02^{n} = 100$ geeft $n = \ln 100/\ln 1.02 = 233$.
\end{solution}

\begin{solution}{exo:b3:sensory-systems:6}
$P = 1 - e^{-a}\sum_{k=0}^{5}a^{k}/k!$: $a = 2$: $0.017$; $a = 6$:
$0.55$; $a = 12$: $0.98$. De helft van de flitsen wordt gezien bij
$a \approx 5.7$; bij een opname van \qty{10}{\%} moet de flits ongeveer
$57$ fotonen aan het hoornvlies leveren.
\end{solution}

\begin{solution}{exo:b3:sensory-systems:7}
$55/3\times 1.3 \approx 24$; in \omterm{def:b3:sensory-systems:cochlea}{decibel} van druk, $20\log_{10}24
\approx \qty{28}{dB}$. Vergroeide beentjes kunnen de trilling niet
doorgeven, zodat het geluid het \omterm{def:b3:sensory-systems:cochlea}{slakkenhuis} alleen door geleiding via
de schedel bereikt, zonder de winst van het middenoor: een
geleidingsverlies van zo'n \qtyrange{30}{50}{dB}.
\end{solution}

\begin{solution}{exo:b3:sensory-systems:8}
Het kanaal van de \omterm{def:b3:sensory-systems:cochlea}{haarcel} wordt rechtstreeks geopend doordat de
\omterm{def:b3:sensory-systems:cochlea}{tipverbinding} eraan trekt --- geen boodschapper, geen enzym --- en
antwoordt dus in microseconden; de fotoreceptor versterkt één foton via
twee enzymatische trappen, en dat kost tientallen milliseconden. Het
gehoor wint de nauwkeurige tijdsordening die nodig is om geluiden aan
vertragingen tussen de oren te lokaliseren en de fase te volgen; het
zien wint de gevoeligheid om afzonderlijke fotonen te tellen, en
betaalt in snelheid.
\end{solution}

\begin{solution}{exo:b3:sensory-systems:9}
$H(0.05) = -0.05\ln 0.05 - 0.95\ln 0.95 = 0.150 + 0.049 = 0.199$.
Mens: $\ln\binom{400}{20} \approx 400\times 0.199 = 80$, ongeveer
$10^{34}$ patronen. Muis: $\ln\binom{1100}{55} \approx 1100\times 0.199
= 219$, ongeveer $10^{95}$.
\end{solution}

\begin{solution}{exo:b3:sensory-systems:10}
Gemiddeld $500\times 0.1/40 = 1.25$ spontane gebeurtenissen per
venster. $P(k
\le 5) = e^{-1.25}(1 + 1.25 + 0.78 + 0.33 + 0.10 + 0.03) = 0.998$, dus
$P(k \ge 6) \approx 2\times 10^{-3}$. Bij een drempel van één zou er in
de meeste vensters een spontane gebeurtenis optreden
($1 - e^{-1.25} = 0.71$) en zou het donker vol flitsen zijn; bij zes
komen valse meldingen eens per vijfhonderd vensters. De drempel wordt
gelegd waar de ruis van de \omterm{def:b3:sensory-systems:retina}{staafjes} hem oplegt, niet hun gevoeligheid.
\end{solution}

\begin{solution}{exo:b3:sensory-systems:11}
Grootste vertraging $0.2/340 = \qty{590}{\micro s}$ (geluid van één
zijde). Nabij het midden is $\Delta t \approx (d/c)\sin\theta$, dus
komt $\qty{10}{\micro
s}$ overeen met $\sin\theta = 0.017$, ongeveer \qty{1}{\degree}. Boven
\qty{4}{kHz} is de periode (\qty{250}{\micro s}) korter dan waarop de
neuronen kunnen vergrendelen, en overspant een vertraging van enkele
honderden microseconden meer dan één cyclus, wat de fase dubbelzinnig
maakt; het brein gebruikt dan het niveauverschil tussen de oren en de
plaatscode.
\end{solution}

\begin{solution}{exo:b3:sensory-systems:12}
Ver van de rand geldt $r = I(1 - 2w)$: $0.6$ aan de donkere kant, $1.2$
aan de heldere. Bij de laatste donkere receptor ($I = 1$, buren $1$ en
$2$): $r
= 1 - 0.2\times 3 = 0.4$; bij de eerste heldere: $r = 2 - 0.6 = 1.4$.
De sprong wordt overdreven tot een doorschot naar beneden en naar boven
--- de banden van Mach. Het \omterm{def:b3:sensory-systems:retina}{netvlies} wint contrast en randen, en een
kleiner dynamisch bereik om over te zenden; het verliest de trouw aan
absolute niveaus, en brengt begoochelingen van helderheid voort.
\end{solution}

\begin{solution}{pb:b3:sensory-systems:1}
\textbf{1.} $0.01/3.6\times 10^{-19} = 2.8\times 10^{16}$ fotonen per
seconde.
\textbf{2.} Pupiloppervlak $\pi(3.5\times 10^{-3})^{2} = 3.8\times 10^{-5}$
m$^{2}$. Op \qty{1}{km}: aandeel $3.8\times 10^{-5}/1.26\times 10^{7} =
3\times 10^{-12}$, $8.5\times 10^{4}$ fotonen per seconde. Op
\qty{10}{km}: $850$ per seconde.
\textbf{3.} In \qty{0.1}{s} met \qty{10}{\%} opname: $850$ op
\qty{1}{km}, $8.5$ op \qty{10}{km} --- beide boven zes; de kaars wordt
op tien kilometer gezien (meestal).
\textbf{4.} Zes opgenomen per venster vergt $600$ fotonen per seconde in
de pupil: $4\pi d^{2} = 3.8\times 10^{-5}\times 2.8\times 10^{16}/600
= 1.8\times 10^{9}$ m$^{2}$, $d \approx \qty{12}{km}$. Daar is $P(k \ge
6 \mid a = 6) = 0.55$: de kaars wordt in ongeveer de helft van de
vensters gezien.
\textbf{5.} $e^{-6} = 0.0025$. Aan de drempel schommelt het aantal van
venster tot venster --- $6 \pm 2.4$ --- zodat de bron lijkt te komen en
te gaan: het fonkelen van een zwakke ster (waaraan de dampkring nog het
hare toevoegt).
\textbf{6.} De schommeling van de achtergrond, $\sqrt{10^{4}} = 100$
fotonen per venster, overspoelt de $8.5$ van de kaars: waarneming vergt
dat het signaal de ruis van de achtergrond overtreft, niet de drempel
van het donkere oog.
\textbf{7.} $\qty{1}{pA}/200 = \qty{5}{fA}$ per kanaal; $5\times
10^{-15}/1.6\times 10^{-19} = 3\times 10^{4}$ ionen per seconde.
\textbf{8.} Een dertigste; ongeveer dertig gelijktijdige fotonen sluiten
elk kanaal en verzadigen het \omterm{def:b3:sensory-systems:retina}{staafje}.
\textbf{9.} $10^{-12}\times 0.2 = 2\times 10^{-13}$ C, $1.3\times 10^{6}$
kationen: één foton beheerst ruim een miljoen ladingen.
\textbf{10.} Het \omterm{def:b3:sensory-systems:retina}{staafje} verzadigt binnen milliseconden, met elk kanaal
gesloten, en zijn \omterm{def:b3:sensory-systems:retina}{rodopsine} wordt sneller gebleekt dan hersteld;
\omterm{def:b3:sensory-systems:retina}{kegeltjes}, met een kleinere versterking en een snellere uitschakeling,
blijven antwoorden en verschuiven hun bereik over zes decaden licht.
\textbf{11.} Een hoge versterking vergt een lange integratie (elke trap
kost tijd en het antwoord wordt gerekt om op te tellen), dus is het
\omterm{def:b3:sensory-systems:retina}{staafje} traag en gevoelig; de kleinere versterking en snellere
uitschakeling van het \omterm{def:b3:sensory-systems:retina}{kegeltje} geven een antwoord in \qty{20}{ms}, snel
genoeg voor beweging, ten koste van honderd benodigde fotonen.
\textbf{12.} $10^{8}/40 = 2.5\times 10^{6}$ spontane isomerisaties per
seconde over het hele \omterm{def:b3:sensory-systems:retina}{netvlies} --- maar slechts $1.25$ per venster in
elke plek van $500$ \omterm{def:b3:sensory-systems:retina}{staafjes}, ver onder de drempel van zes, zodat de
ruis plek voor plek wordt verworpen; de drempel en het samennemen
bestaan juist daarvoor.
\textbf{13.} \qty{0.1}{W/m^2}; op \qty{5.5e-5}{m^2}, \qty{5.5e-6}{W}.
\textbf{14.} $5.5\times 10^{-6}\times 7200 = \qty{0.04}{J}$ --- ongeveer
vierhonderd regendruppels in twee uur.
\textbf{15.} $10^{-12}\times 5.5\times 10^{-5} = 5.5\times 10^{-17}$ W;
in \qty{10}{ms}, $5.5\times 10^{-19}$ J --- ongeveer de energie van één
of twee zichtbare fotonen.
\textbf{16.} $0.3/5000 = 6\times 10^{-5}$ rad, $0.003^{\circ}$; de
uitwijking is ongeveer drie waterstofatomen.
\textbf{17.} $10^{(110 - 85)/10} = 316$: twee uur bij \qty{110}{dB}
leveren de energie van $630$ uur bij \qty{85}{dB} --- bijna een maand
werkdagen.
\textbf{18.} Ongeveer $120$ net merkbare stappen, één per \omterm{def:b3:sensory-systems:cochlea}{decibel}.
\textbf{19.} $2^{400} = 10^{400\log_{10}2} = 10^{120}$.
\textbf{20.} $\ln\binom{400}{20} \approx 400\,H(0.05) \approx 80$, min
een correctie van Stirling van ongeveer $2$: zo'n $10^{34}$ patronen.
\textbf{21.} Tel de receptoren die in het ene patroon zitten en in het
andere niet: $5 +
5 = 10$ van de $40$ (een afstand van Hamming); patronen die voor drie
kwart overlappen sturen vrijwel dezelfde glomeruli aan en worden als
vrijwel dezelfde geur gelezen.
\textbf{22.} $10^{12}/10^{34} = 10^{-22}$ van de patronen. De code is
niet de grens: receptoren zijn ruizig, vele zijn gelijk afgestemd zodat
hun activiteiten samenhangen, echte moleculen brengen slechts een
kleine deelverzameling van de patronen voort, en het onderscheidings-
en geheugenvermogen van het brein voor patronen is eindig.
\textbf{23.} $\ln\binom{1100}{20} \approx 1100\,H(0.018) = 1100\times
0.091 \approx 100$, ongeveer $10^{43}$ --- zo'n $10^{9}$ maal het
menselijke aantal patronen van twintig receptoren (en veel meer nog als
de muis meer receptoren per geur gebruikt).
\textbf{24.} De reuk in het algemeen verandert nauwelijks --- de code is
redundant en andere receptoren dragen elke geurstof --- maar een
geurstof die bij lage concentratie van die receptor afhangt wordt
onwaarneembaar of verandert van karakter: een specifieke anosmie, zoals
bij de vele mensen die androstenon niet kunnen ruiken.
\textbf{25.} De kaars wordt tot ongeveer \qty{12}{km} gezien;
drempelgeluid levert in \qty{10}{ms} $5.5\times 10^{-19}$ J aan het
trommelvlies, één of twee fotonen; een menselijke neus kan zo'n
$10^{34}$ patronen van twintig receptoren vormen.
\end{solution}
